% Ad Familiares 10.12 (ad-familiares-10-12)
% Source: https://raw.githubusercontent.com/PerseusDL/canonical-latinLit/master/data/phi0474/phi056/phi0474.phi056.perseus-lat1.xml
% Fetched by scripts/fetch_latin.py

% Dateline (Perseus): Scr. Romae iii Id. Apr. a. 711 (43).

\ciceroLetterOpener{CICERO PLANCO.}

\ciceroSection{1} etsi rei p. causa maxime gaudere debeo tantum ei te praesidi, tantum opis attulisse extremis paene temporibus, tamen ita te victorem complectar re p. reciperata, ut magnam partem mihi laetitiae tua dignitas adfert, quam et esse iam et futuram amplissimam intellego. cave enim putes ullis umquam litteras gratiores quam tuas in senatu esse recitatas ; idque contigit cum meritorum tuorum in rem p. eximia quadam magnitudine tum verborum sententiarumque gravitate. quod mihi quidem minime novum, qui et te nossem et tuarum litterarum ad me missarum promissa meminissem et haberem a Furnio nostro tua penitus consilia cognita, sed senatui maiora visa sunt quam erant exspectata, non quo umquam de tua voluntate dubitasset, sed nec quantum facere posses nec quo progredi velles exploratum satis habebat.

\ciceroSection{2} itaque cum a. d. vii Idus Aprilis mane mihi tuas litteras M. Varisidius reddidisset easque legissem, incredibili gaudio sum elatus ; cumque magna multitudo optimorum virorum et civium me de domo deduceret, feci continuo omnis participes meae voluptatis. interim ad me venit Munatius noster, ut consuerat. at ego ei litteras tuas, nihildum enim sciebat; nam ad me primum Varisidius, idque sibi a te mandatum esse dicebat. Paulo post idem mihi Munatius eas litteras legendas dedit quas ipsi miseras, et eas quas publice.

\ciceroSection{3} placuit nobis ut statim ad Cornutum, pr. urb., litteras deferremus, qui, quod consules aberant, consulare munus sustinebat more maiorum. senatus est continuo convocatus frequensque convenit propter famam atque exspectationem tuarum litterarum. recitatis litteris oblata religio Cornuto est pullariorum admonitu non satis diligenter eum auspiciis operam dedisse, idque a nostro conlegio comprobatum est. itaque res dilata est in posterum. eo autem die magna mihi pro tua dignitate contentio cum Servilio ; qui cum gratia effecisset ut sua sententia prima pronuntiaretur, frequens eum senatus reliquit et in alia omnia discessit, meaeque sententiae, quae secunda pronuntiata erat, cum frequenter adsentiretur senatus, rogatu . Servili P. Titius intercessit. res in posterum dilata.

\ciceroSection{4} venit paratus Servilius Iovi ipsi iniquus, cuius in templo res agebatur. hunc quem ad modum fregerim quantaque contentione Titium intercessorem abiecerim, ex aliorum te litteris malo cognoscere ; unum hoc ex meis : senatus gravior, constantior, amicior tuis laudibus esse non potuit quam tum fuit, nec vero tibi senatus amicior quam cuncta civitas ; mirabiliter enim populus R. universus et omnium generum ordinumque consensus ad liberandam rem p. conspiravit.

\ciceroSection{5} Perge igitur, ut agis, nomenque tuum commenda immortalitati atque haec omnia, quae habent speciem gloriae conlectam inanissimis splendoris insignibus contemne, brevia, fucata, caduca existima. verum decus in virtute positum est, quae maxime inlustratur magnis in rem p. meritis. eam facultatem habes maximam. quam quoniam complexus es et tenes, perfice ut ne minus res p. tibi quam tu rei p. debeas. me tuae dignitatis non modo fautorem sed etiam amplificatorem cognosces. id cum rei p. quae mihi vita est mea carior, tum nostrae necessitudini debere me iudico. atque in his curis, quas contuli ad dignitatem tuam, cepi magnam voluptatem, quod bene cognitam mihi T. Munati prudentiam et fidem magis etiam perspexi in eius incredibili erga te benevolentia et diligentia. iii Idus Apr.
